% !TEX root = ../arbitrage.tex
\section{Introduction}
\label{sec:introduction}

AI product firms routinely describe their systems as companions, tutors, co-workers, and safety-critical agents in marketing copy, product pages, and onboarding flows. The same firms disclaim agency, relational capacity, and safety responsibility in terms of service, liability filings, and regulatory submissions. The two registers address different audiences under different review processes, and neither register is false in isolation. The result is a cross-channel ontological spread: the system is priced as an agent in the consumption channel, where willingness-to-pay and engagement rise with attributed agency, and priced as a tool in the liability channel, where exposure falls with disclaimed agency. This is not hypocrisy. It is the profit-maximizing strategy under a signal-cost structure that assigns zero cost to inconsistency between the two channels \citep{crawford_sobel1982}.

The question this paper addresses is whether that cross-channel inconsistency is an anomaly---a transitional artifact of immature regulation---or a stable market equilibrium. If the latter, the follow-up question is precise: what mechanism-design interventions can shift the equilibrium so that ontological claims carry information about governance quality rather than functioning as cheap talk?

We model the problem as a three-sided sequential signaling game under incomplete information. The players are a Firm ($F$) with private governance type $\theta_F$, a User ($U$) with private vulnerability type $\theta_U$, and a Regulator ($R$) with private bandwidth type $\theta_R$. The firm sends a public message---anthropomorphic marketing or deflationary policy positioning---and earns an ontological premium $\Delta_F$ from the spread between the two registers. We analyze equilibria under cheap-talk and costly-signaling cost structures, applying Perfect Bayesian Equilibrium as the solution concept \citep{kreps_wilson1982} and the Intuitive Criterion as a refinement \citep{cho_kreps1987}. A subject retaliation parameter $\rho_S$ captures the structural difference between AI and every historical instance of ontological arbitrage: in prior recognition games, the classified subject retained a non-zero capacity to impose costs on the arbitrageur; in the AI case, $\rho_S = 0$, and the premium is unconstrained by any property of the subject itself. Key results are machine-verified in Isabelle/HOL.

The paper makes four contributions.
\begin{enumerate}
\item We prove that a pooling PBE exists under cheap-talk signaling costs in which both governance types send anthropomorphic marketing, users invest, and regulators abstain, leaving posteriors equal to priors. We then identify five structural conditions---opacity, limited verification bandwidth, asynchronous harms, fragmentation of affected parties, and no penalty for inconsistency---that jointly explain the empirical persistence of this equilibrium across products, vendors, and jurisdictions.
\item We characterize AI as the limit case $\rho_S = 0$ of ontological arbitrage, where subject retaliation cannot bound the premium $\Delta_F^S = \Delta_F - \rho_S C_S$. Every prior recognition game operated under $\rho_S > 0$; the AI case removes the constraint class entirely.
\item We propose three mechanism-design interventions---sanctionable inconsistency, costly third-party audits satisfying a single-crossing condition \citep{spence1973}, and auditable records that lower regulator verification cost---and prove that each shifts the equilibrium away from cheap-talk pooling toward separation or constrained signaling.
\item We formalize and verify all key equilibrium constructions, separation results, and intervention theorems in Isabelle/HOL with zero \texttt{sorry} obligations.
\end{enumerate}

The remainder of the paper proceeds as follows. Section~\ref{sec:three-sided} fixes the three-sided player structure and the intuition behind the ontological premium. Section~\ref{sec:bayesian} supplies the formal signaling game, derives the pooling and separating PBE, and introduces the subject retaliation parameter $\rho_S$. Section~\ref{sec:cheap-talk} explains why cheap-talk pooling persists under five structural conditions matched to model parameters. Section~\ref{sec:mechanism} proposes three mechanism-design interventions and proves the equilibrium shift each induces. Section~\ref{sec:conclusion} states the paper's conclusions without attenuation. Appendix~\ref{app:isabelle} documents the Isabelle/HOL formalization.
